\chapter{Положительная толщина, исключённый объём и барьер пересоединения}
\label{ch:v6-single-thread-excluded-volume-reconnection-barrier-gate}

\section{Проверяемое утверждение}

Предыдущая глава построила сбалансированный Real-цикл, но оставила
открытым различие между абстрактным маршрутом и физической нитью. В этой
главе проверяется буквальная ранняя гипотеза: одна нить имеет положительное
сечение, её несоседние участки не могут занимать одно место, а изменение
типа узла требует бесконечного или родительски большого барьера.

Строгий ledger уже содержит локальное натяжение, производные и кривизну
полей, внутренние барьеры потери ранга, winding- и хопфовы метки. Ни один
из этих результатов нельзя заранее переименовать в исключённый объём.

\section{Положительная толщина как условие на reach}

Для гладкой замкнутой осевой кривой $X:S^1\to M$ трубка радиуса $a$ не
самопересекается только при
\begin{equation}
 a<\operatorname{reach}(X).
\end{equation}
Схематически reach ограничивается одновременно локальным радиусом
кривизны и половиной расстояния между несоседними участками. Поэтому одна
только конечная кривизна не исключает почти касательного самоконтакта.
Именно это различие лежит в основе геометрии толщины и ropelength
\cite{single-thread-gonzalez-maddocks1999,
single-thread-cantarella-fu-kusner-sullivan2014}.

Текущая ширина эффективного вихря $0.57328995$ не является радиусом
фундаментальной нити: она получена после континуального усреднения и не
фиксирует микроскопический cutoff.

\section{Условная упаковочная оценка}

Пусть $K$ непересекающихся круглых сечений радиуса $a$ одновременно
проходят через грубую квадратную ячейку линейного размера $\ell$. Даже
при предельной гексагональной упаковке необходимо
\begin{equation}
 K\pi a^2\leq \frac{\pi}{2\sqrt3}\ell^2,
 \qquad
 \frac a\ell\leq\frac{1}{\sqrt{2\sqrt3K}}.
 \label{eq:v6-single-thread-packing-bound}
\end{equation}

Для максимальной рабочей кратности $K=22696855$ это даёт
\begin{equation}
 \frac a\ell\leq1.1278\cdot10^{-4}.
\end{equation}
Для всех $50370998$ Real-ориентированных посещений условная граница равна
$7.5703\cdot10^{-5}$.

Эта оценка не доказывает физическое число витков. Напротив, она показывает,
что целочисленные кратности предыдущей главы нельзя одновременно считать
буквальным количеством планковских трубок в одной планковской ячейке.
Они были построены как рациональные аппроксиманты моментов $Q_*$.

Сохраняется только более слабая возможность: одна нить образует конечную
cutoff-ткань, а грубая ячейка содержит много более мелких микроскопических
ячеек. Ни отношение масштабов, ни точное целочисленное слово родителем ещё
не выведены.

\section{Локальные энергии не видят самоконтакт}

Рассмотрим два прямых единичных участка на расстоянии $\delta$. Их полная
длина равна двум, а локальная кривизна равна нулю при любом $\delta>0$.
Следовательно, натяжение Намбу--Гото и локальный изгиб остаются конечными
при $\delta\to0$.

Для сравнения нелокальный контактный функционал с ядром второго порядка
даёт
\begin{equation}
 I_2(\delta)
 =\int_0^1\!\!\int_0^1
 \frac{ds\,dt}{(s-t)^2+\delta^2}
 =\frac{2}{\delta}\arctan\frac1\delta
 -\log\frac{1+\delta^2}{\delta^2},
\end{equation}
и
\begin{equation}
 \delta I_2(\delta)\longrightarrow\pi.
\end{equation}
Такой член действительно создаёт бесконечную цену контакта. Родственные
самоотталкивающиеся функционалы были введены в энергии узлов О'Хары и
Мёбиуса \cite{single-thread-ohara1991,single-thread-freedman-he-wang1994}.

Но в текущем родителе нет двойного интеграла по двум независимым точкам
одной пространственной нити, hard-core ограничения или выведенного
тангенциально-точечного ядра. Добавление $I_2$ с новым коэффициентом было
бы новой микроскопической аксиомой.

\section{Топологический заряд не запрещает пересоединение}

Winding, хопфов заряд и остаточная голономия классифицируют конфигурации,
пока поле остаётся в соответствующем вакуумном многообразии. В ядре
дефекта амплитуда может покинуть вакуумную орбиту. Поэтому столкновение
двух конечных ядер может пройти через конечную конфигурацию, в которой
локальная топологическая метка не определена.

Это не только логическая возможность: полевые вихревые струны модели
Абеля--Хиггса классически пересоединяются при столкновении
\cite{single-thread-hanany-hashimoto2005,
single-thread-achucarro-verbiest2010}. Напротив, устойчивые хопфовы узлы
в моделях Фаддеева--Скирма удерживаются полным полевым действием и
хопфовым сектором, а не постулатом материальной неразрывной верёвки
\cite{single-thread-faddeev-niemi1997}.

\begin{theorem}[Граница текущего родителя]
Существующие локальные натяжение, кривизна, ранговые барьеры и
топологические метки не выводят положительный reach фундаментальной нити
и не создают бесконечный барьер пересоединения. Буквальная неразрывная
одна нить текущим действием не получена.
\end{theorem}

\section{Что сохраняется после отрицательного результата}

Не закрыты:
\begin{enumerate}
 \item эффективное чтение $Q,T,B$ как моментов многопроходной ткани;
 \item конечная cutoff-укладка при отдельно выведенной иерархии
 масштабов;
 \item устойчивый хопфов солитон локального полевого действия;
 \item коллективная инерция уже построенного дефекта.
\end{enumerate}

Закрыто более сильное утверждение, будто имеющийся родитель уже заставляет
все проходы быть одной физически неразрывной нитью планковской толщины.

Следующий гейт должен принять архитектурное решение между двумя ветвями:
\begin{enumerate}
 \item вывести микроскопический масштаб $a/\ell$ и hard-core правило из
 уже существующего cutoff/спектрального носителя;
 \item отказаться от буквальной верёвки и продолжить с устойчивыми
 хопфовыми солитонами эффективной ткани.
\end{enumerate}

Нулевая гипотеза: существующие cutoff, конечный носитель или спектральная
поляризация уже задают дискретную иерархию $a\ll\ell$ и запрет изменения
связности без нового параметра.

Стоп-критерий: если отношение масштабов и hard-core правило отсутствуют,
буквальная одна нить переводится в историческую метафору, а рабочей
физической ветвью остаётся хопфов солитон.

\begin{protocol}
\begin{enumerate}
 \item фундаментальный радиус нити: \textbf{не выведен};
 \item положительный reach: \textbf{не выведен};
 \item равенство толщины нити и масштаба грубой ячейки:
 \textbf{несовместимо с условной многопроходной упаковкой};
 \item натяжение расходится при контакте: \textbf{нет};
 \item локальная кривизна расходится при параллельном контакте:
 \textbf{нет};
 \item нелокальная контактная энергия могла бы расходиться:
 \textbf{да, но она отсутствует};
 \item топология запрещает пересоединение конечных ядер: \textbf{нет};
 \item буквальная неразрывная нить выведена: \textbf{нет};
 \item эффективная многопроходная ткань закрыта: \textbf{нет};
 \item следующий тест: \textbf{решение ветви и иерархии масштабов}.
\end{enumerate}
\end{protocol}

Машинный сертификат сохранён в
\path{s2t_v6_single_thread_excluded_volume_reconnection_barrier_gate_results.json}.

\begin{thebibliography}{9}
\bibitem{single-thread-gonzalez-maddocks1999}
O.~Gonzalez, J.~H.~Maddocks,
\emph{Global Curvature, Thickness, and the Ideal Shapes of Knots},
PNAS 96 (1999), 4769--4773.

\bibitem{single-thread-cantarella-fu-kusner-sullivan2014}
J.~Cantarella, J.~H.~G.~Fu, R.~Kusner, J.~M.~Sullivan,
\emph{Ropelength Criticality},
Geometry \& Topology 18 (2014), 1973--2043.

\bibitem{single-thread-ohara1991}
J.~O'Hara,
\emph{Energy of a Knot},
Topology 30 (1991), 241--247.

\bibitem{single-thread-freedman-he-wang1994}
M.~H.~Freedman, Z.-X.~He, Z.~Wang,
\emph{Möbius Energy of Knots and Unknots},
Annals of Mathematics 139 (1994), 1--50.

\bibitem{single-thread-hanany-hashimoto2005}
A.~Hanany, K.~Hashimoto,
\emph{Reconnection of Colliding Cosmic Strings},
JHEP 06 (2005), 021; arXiv:hep-th/0501031.

\bibitem{single-thread-achucarro-verbiest2010}
A.~Achúcarro, G.~J.~Verbiest,
\emph{Higher Order Intercommutations in Cosmic String Collisions},
Physical Review Letters 105 (2010), 021601.

\bibitem{single-thread-faddeev-niemi1997}
L.~Faddeev, A.~J.~Niemi,
\emph{Stable Knot-Like Structures in Classical Field Theory},
Nature 387 (1997), 58--61.
\end{thebibliography}